% =============================================================================
% Implementação de Baixo Nível da Biblioteca CKKS — versão resumida (5–6 pág.)
% Compilar: pdflatex dissertacao_baixo_nivel_resumido.tex  (2 passagens)
% =============================================================================
\documentclass[11pt,a4paper]{article}

\usepackage[utf8]{inputenc}
\usepackage[T1]{fontenc}
\usepackage[brazil]{babel}

\usepackage{amsmath}
\usepackage{amssymb}
\usepackage{amsthm}

\usepackage{tikz}
\usetikzlibrary{arrows.meta, positioning, calc, decorations.pathreplacing,
                shapes.geometric}

\usepackage{booktabs}
\usepackage[a4paper, top=2.5cm, bottom=2.5cm, left=2.8cm, right=2.8cm]{geometry}

\usepackage[colorlinks=true, allcolors=blue,
            pdfauthor={Yuri Dimitre Dias de Faria},
            pdftitle={Implementação de Baixo Nível da Biblioteca CKKS -- resumo}]{hyperref}
\usepackage[brazilian, nameinlink]{cleveref}

\newtheorem{definition}{Definição}

\newcommand{\Rq}{\mathcal{R}_q}

% Títulos das seções em português
\renewcommand\abstractname{Resumo}

% =============================================================================
\begin{document}

\title{\textbf{Implementação de Baixo Nível da Biblioteca CKKS\\
para Aprendizado Federado Preservador de Privacidade}\\[0.4em]
\large Síntese técnica}
\author{Yuri Dimitre Dias de Faria\\
\small Instituto de Computação -- UNICAMP\\
\small Orientador: Prof. Dr. Allan Mariano de Souza}
\date{Junho de 2026}
\maketitle

% =============================================================================
\section{Introdução}
\label{sec:intro}
% =============================================================================

Este documento sintetiza a arquitetura de baixo nível do \emph{backend} C++ da
biblioteca CKKS customizada, desenvolvida para suportar operações de cifragem
homomórfica em cenários de Aprendizado Federado (FL). O material pressupõe
familiaridade com o pano de fundo do esquema CKKS, com o modelo de erro da
agregação homomórfica e com o problema RLWE, cobertos no Capítulo~4 da dissertação
associada~\cite{Cheon2017,sealmanual}; o foco aqui é a tradução dessas estruturas
matemáticas em código executável de alto desempenho.

O núcleo da biblioteca, denominado \texttt{ckks\_core}, é implementado em C++17 e
organizado em seis módulos hierárquicos (\cref{tab:modulos}). Todas as unidades de
tradução são compiladas para uma biblioteca estática com as diretivas
\texttt{-O3 -march=native -ffast-math -fopenmp-simd}.

\begin{table}[ht]
\centering
\caption{Módulos do núcleo \texttt{ckks\_core}.}
\label{tab:modulos}
\small
\begin{tabular}{lll}
\toprule
\textbf{Módulo} & \textbf{Arquivos principais} & \textbf{Responsabilidade} \\
\midrule
1 – Aritmética & \texttt{math\_utils.h/.cpp} & \texttt{mod\_add/sub}, Montgomery, primos NTT \\
2 – Polinômio  & \texttt{poly.h/.cpp}        & Tipo \texttt{Poly}, adição vetorizada \\
3 – NTT        & \texttt{ntt.h/.cpp}         & Transformada negacíclica $O(N\log N)$ \\
4 – Amostragem & \texttt{sampler.h/.cpp}     & Gaussiana discreta, distribuições ternárias \\
5 – Chaves/Cifragem & \texttt{keygen.h}, \texttt{ckks\_core.cpp} & Geração de chaves, \texttt{encrypt}, \texttt{decrypt} \\
6 – Avaliação  & \texttt{evaluator.h/.cpp}, \texttt{encoder.h/.cpp}, \texttt{serialization.cpp} & Adição, codificação, serialização \\
\bottomrule
\end{tabular}
\end{table}

Em contraste com implementações de propósito geral como Microsoft
SEAL~\cite{sealmanual}, o \texttt{ckks\_core} adota duas simplificações intencionais,
motivadas pelo foco exclusivo em agregação federada: (i)~utiliza um \textbf{único primo
NTT} $q \equiv 1 \pmod{2N}$, com $q < 2^{62}$, em lugar da cadeia RNS
$Q = \prod_j q_j$ do modelo completo discutido no Capítulo~4; (ii)~opera em
\textbf{regime somente-aditivo}, sem \emph{rescaling}, relinearização, chaves de Galois
nem \emph{bootstrapping}. O orçamento de ruído, portanto, não decresce ao longo das
rodadas de agregação.

% =============================================================================
\section{Aritmética Modular e Redução de Montgomery}
\label{sec:arithmetic}
% =============================================================================

\paragraph{Anel de polinômios.}
A biblioteca opera sobre o anel quociente
\begin{equation}
  \Rq = \mathbb{Z}_q[X]\,/\,(X^N+1),
  \label{eq:ring}
\end{equation}
onde $q$ é um primo e $N$ é uma potência de dois. A relação $X^N \equiv -1
\pmod{X^N+1}$ confere ao anel estrutura negacíclica: o produto de polinômios é
reduzido módulo $X^N+1$ implicitamente pela NTT, descrita na \cref{sec:ntt}.

\paragraph{Aritmética sem desvios condicionais.}
Adição e subtração modulares são implementadas sem desvios (\emph{branchless})
por meio do truque da máscara binária. Para $a, b \in [0,q)$:
\[
  \texttt{mod\_add}(a,b,q) = (a+b) - \bigl(q \;\mathbin{\&}\; -(a+b \geq q)\bigr),
\]
onde a negação do resultado booleano em 64 bits produz a máscara $0$ ou
$2^{64}-1$. A subtração \texttt{mod\_sub} usa raciocínio simétrico. Ambas são
declaradas \texttt{inline noexcept} e mapeiam diretamente a instruções AVX2
(\texttt{vpsubq}/\texttt{vpcmpgtq}/\texttt{vpand}) sob \texttt{-march=native}.

\paragraph{Redução de Montgomery.}
Multiplicações modulares no caminho crítico (borboletas NTT, cifragem) empregam
a redução de Montgomery~\cite{Montgomery1985} com $R = 2^{64}$. Cada elemento
$a \in \mathbb{Z}_q$ é representado na \emph{forma de Montgomery}
$\hat{a} = a \cdot R \bmod q$; a multiplicação
$\widehat{a \cdot b} = \texttt{REDC}(\hat{a} \cdot \hat{b})$ nunca realiza divisão
por $q$, substituindo-a por um deslocamento de 64 bits e uma adição. As três
constantes do contexto \texttt{MontgomeryCtx} ($q_{\text{inv}} = -q^{-1} \bmod 2^{64}$,
$r_2 = R^2 \bmod q$ e $q$) são pré-computadas uma única vez por par $(N, q)$.
O valor $q_{\text{inv}}$ é obtido por cinco iterações de Newton $x \leftarrow x(2-qx)$
sobre $\mathbb{Z}/2^{64}\mathbb{Z}$; o valor $r_2$ via redução de 128 bits. A
\cref{fig:montgomery_flow} ilustra o padrão de uso: as conversões ocorrem uma única
vez para cada conjunto de operandos, enquanto as multiplicações intermediárias
permanecem inteiramente no domínio de Montgomery.

\begin{figure}[ht]
\centering
\begin{tikzpicture}[
  node distance = 0.5cm and 1.3cm,
  box/.style  = {draw, rounded corners=3pt, minimum width=2.5cm, minimum height=0.65cm,
                 align=center, font=\small},
  io/.style   = {draw, ellipse, minimum width=1.8cm, minimum height=0.65cm,
                 align=center, font=\small},
  arr/.style  = {-{Stealth[length=5pt]}, thick}
]
  \node[io]  (a)  {$a \in \mathbb{Z}_q$};
  \node[io,  below=0.7cm of a] (b) {$b \in \mathbb{Z}_q$};
  \node[box, right=1.2cm of a]  (am)  {$\hat{a}=aR\bmod q$};
  \node[box, right=1.2cm of b]  (bm)  {$\hat{b}=bR\bmod q$};
  \node[box, right=1.2cm of am, yshift=-0.67cm] (mul) {\texttt{mont\_mul}\\$\hat{a}\cdot\hat{b}$};
  \node[box, right=1.2cm of mul] (res) {$\widehat{ab}$};
  \node[io,  right=1.1cm of res] (out) {$ab\in\mathbb{Z}_q$};
  \draw[arr] (a)   -- node[above,font=\footnotesize]{\texttt{to\_mont}}   (am);
  \draw[arr] (b)   -- node[below,font=\footnotesize]{\texttt{to\_mont}}   (bm);
  \draw[arr] (am) |- (mul);
  \draw[arr] (bm) |- (mul);
  \draw[arr] (mul) -- node[above,font=\footnotesize]{REDC}                (res);
  \draw[arr] (res) -- node[above,font=\footnotesize]{\texttt{from\_mont}} (out);
\end{tikzpicture}
\caption{Fluxo de uso do domínio de Montgomery. Conversões ocorrem uma única vez;
multiplicações intermediárias usam \texttt{mont\_mul} sem divisão por $q$.}
\label{fig:montgomery_flow}
\end{figure}

\paragraph{Geração de primos NTT-amigáveis.}
A função \texttt{generate\_proth\_prime(bit\_size, N)} encontra o maior primo da forma
$q = 2Nk+1$ com exatamente \texttt{bit\_size} bits ($\leq 62$), garantindo
$q \equiv 1 \pmod{2N}$ e $q < 2^{62}$. A busca é descendente em passos de $2N$; cada
candidato é submetido ao teste de Miller--Rabin determinístico com 12 testemunhas fixas,
suficientes para cobrir todo o intervalo $[0, 2^{62})$ sem falsos
positivos~\cite{Pomerance1980}.

% =============================================================================
\section{Transformada Numérica Teórica Negacíclica}
\label{sec:ntt}
% =============================================================================

A multiplicação de polinômios em $\Rq$ domina o custo computacional da cifragem e da
decifragem. A Transformada Numérica Teórica (NTT) reduz a complexidade de $O(N^2)$
para $O(N\log N)$~\cite{LongaNaehrig2016}. A variante \emph{negacíclica} avalia os
polinômios nas $N$ raízes primitivas $2N$-ésimas da unidade
$\psi, \psi^3, \ldots, \psi^{2N-1}$, onde $\psi^N \equiv -1 \pmod{q}$, tornando a
redução módulo $X^N+1$ implícita ao algoritmo.

O contexto \texttt{NttCtx} é construído uma única vez por par $(N, q)$ e armazena:
(i)~a raiz $\psi = g^{(q-1)/(2N)} \bmod q$, derivada da menor raiz primitiva $g$ de
$\mathbb{Z}_q^*$; (ii)~duas tabelas de $N$ \emph{twiddle factors}
$\texttt{psi\_pow}[k] = \texttt{to\_mont}(\psi^{\text{bit\_rev}(k)})$ e
$\texttt{psi\_inv\_pow}[k]$, geradas por multiplicação sequencial de Montgomery e
armazenadas em \textbf{ordem bit-reversa} para acesso linear sequencial durante as
borboletas; (iii)~o escalar $N_{\text{inv}}^{\text{mont}} = \texttt{to\_mont}(N^{-1}
\bmod q)$, usado na normalização da NTT inversa.

A NTT direta (Cooley--Tukey DIT) processa $\log_2 N$ estágios com $m$ crescente,
aplicando o fator de giro \emph{antes} de cada soma/diferença. A NTT inversa
(Gentleman--Sande DIF) processa os estágios com $m$ decrescente, aplica o fator
inverso \emph{após} a soma/diferença e normaliza todos os coeficientes por
$N_{\text{inv}}^{\text{mont}}$ ao final. Os dois algoritmos formam um par exato:
nenhuma permutação adicional é necessária após a NTT inversa. Os ponteiros
\texttt{\_\_restrict\_\_} nos laços internos informam ao compilador a ausência de
aliasing, favorecendo a vetorização automática.

A multiplicação polinomial \texttt{poly\_ntt\_mul} aplica a NTT direta a cada operando,
multiplica ponto a ponto usando aritmética de 128 bits (os resíduos de saída da NTT
estão em forma normal, não de Montgomery) e aplica a NTT inversa ao resultado, com
complexidade total $O(N\log N)$.

% =============================================================================
\section{Geração de Chaves, Cifragem e Decifragem}
\label{sec:crypto}
% =============================================================================

\paragraph{Amostragem.}
A classe \texttt{Sampler} encapsula um gerador Mersenne Twister de 64 bits
(\texttt{mt19937\_64}), não copiável. Três distribuições são fornecidas: uniforme em
$[0, q)$; gaussiana discreta $D_{\mathbb{Z},\sigma}$ com rejeição de amostras além de
$6\sigma$ (probabilidade de rejeição $< 2 \times 10^{-9}$ por coeficiente, $\sigma =
3{,}2$); e ternária em $\{-1, 0, 1\}$, com variante de peso de Hamming fixo por
Fisher--Yates parcial.

\paragraph{Geração de chaves.}
O esquema de chaves segue o paradigma RLWE~\cite{LPR2010}. A chave secreta $s$ é um
polinômio ternário com coeficientes em $\{0, 1, q{-}1\}$. A chave pública é o par
$(p_0, p_1) \in \Rq^2$ construído como $p_1 = a$ (uniforme) e $p_0 = e - a \cdot s$,
onde $e$ é um ruído gaussiano pequeno. A invariante resultante
\begin{equation}
  p_0 + p_1 \cdot s = e \pmod{q}
  \label{eq:pk}
\end{equation}
é a base da segurança RLWE: conhecendo apenas $(p_0, p_1)$, recuperar $s$ é
computacionalmente inviável.

\paragraph{Cifragem.}
Dado um texto claro $m \in \Rq$ (um polinômio de coeficientes inteiros, produzido pelo
codificador canônico), a cifragem sorteia um polinômio efêmero $u$ (ternário) e dois
ruídos $e_0, e_1$ (gaussianos) frescos a cada chamada, e calcula
\begin{equation}
  c_1 = p_1 u + e_1, \qquad c_0 = p_0 u + e_0 + m.
  \label{eq:enc}
\end{equation}
As duas multiplicações polinomiais em~\eqref{eq:enc} são realizadas por
\texttt{poly\_ntt\_mul}; as somas são coeficiente a coeficiente com \texttt{mod\_add}.

\paragraph{Decifragem.}
A decifragem computa simplesmente
\begin{equation}
  m' = c_0 + c_1 \cdot s \pmod{q}.
  \label{eq:dec}
\end{equation}
Substituindo~\eqref{eq:enc} em~\eqref{eq:dec} e usando a invariante~\eqref{eq:pk}:
\[
  m' = m + \underbrace{u \cdot e_{\text{chave}} + e_0 + e_1 s}_{\varepsilon},
\]
onde o termo de erro $\varepsilon$ tem coeficientes limitados em norma infinita por
$\approx 2N \cdot 6\sigma \ll q/2$, garantindo recuperação correta de $m$ após
decodificação por $\Delta$. A \cref{fig:pipeline} resume o fluxo completo de cifragem,
agregação e decifragem.

\begin{figure}[ht]
\centering
\begin{tikzpicture}[
  node distance = 0.35cm and 0.75cm,
  box/.style = {draw, rounded corners=3pt, minimum width=2.2cm, minimum height=0.65cm,
                align=center, font=\small},
  io/.style  = {draw, ellipse, minimum width=1.7cm, minimum height=0.65cm,
                align=center, font=\small},
  arr/.style = {-{Stealth[length=5pt]}, thick},
  sub/.style = {font=\scriptsize, text=gray}
]
  \node[io]  (m)   {$m^{(k)}\!\in\!\Rq$};
  \node[box, right=0.65cm of m]   (enc)  {\texttt{encrypt}\\$(p_0,p_1,u,e_0,e_1)$};
  \node[io,  right=0.65cm of enc] (ct)   {$(c_0^{(k)},c_1^{(k)})$};
  \node[box, right=0.65cm of ct]  (add)  {\texttt{add\_inplace}\\$k=1,\ldots,K$};
  \node[io,  right=0.65cm of add] (cta)  {$(c_0',c_1')$};
  \node[box, right=0.65cm of cta] (dec)  {\texttt{decrypt}\\$c_0'+c_1' s$};
  \node[io,  right=0.65cm of dec] (m2)   {$m'\!\approx\!\sum m^{(k)}$};
  \draw[arr] (m)   -- (enc);
  \draw[arr] (enc) -- (ct);
  \draw[arr] (ct)  -- (add);
  \draw[arr] (add) -- (cta);
  \draw[arr] (cta) -- (dec);
  \draw[arr] (dec) -- (m2);
  \node[sub, below=0.1cm of enc] {2$\times$ NTT-mul};
  \node[sub, below=0.1cm of add] {$2N$ \texttt{mod\_add}};
  \node[sub, below=0.1cm of dec] {1$\times$ NTT-mul};
\end{tikzpicture}
\caption{Pipeline de cifragem, agregação homomórfica de $K$ clientes e decifragem.
As operações dominantes são as multiplicações polinomiais via NTT ($O(N\log N)$ cada).}
\label{fig:pipeline}
\end{figure}

% =============================================================================
\section{Avaliação Homomórfica e Codificação Canônica}
\label{sec:eval_enc}
% =============================================================================

\paragraph{Adição homomórfica.}
A classe \texttt{Evaluator} é sem estado além dos parâmetros $(N,q)$. O método
\texttt{add\_inplace} realiza a adição coeficiente a coeficiente de $c_0$ e de $c_1$
de dois textos cifrados com mesmo $(N,q)$, usando dois laços de \texttt{mod\_add}
anotados com \texttt{\#pragma omp simd}, que sob \texttt{-fopenmp-simd} emitem
instruções AVX2 vetoriais sem acionar o \emph{runtime} de \emph{threads} OpenMP. A
correção é imediata pela linearidade:
\[
  (c_0^{(1)}\!+c_0^{(2)}) + (c_1^{(1)}\!+c_1^{(2)})\cdot s
  = (c_0^{(1)}\!+c_1^{(1)}s) + (c_0^{(2)}\!+c_1^{(2)}s) \approx m_1 + m_2.
\]
O ruído acumulado após $K$ adições de textos independentes cresce em norma RMS como
$\sqrt{K}\,\sigma_e$, produzindo erro por \emph{slot}
\begin{equation}
  \varepsilon_K \approx \frac{\sqrt{K}\,\sigma_e}{\Delta},
  \label{eq:noise}
\end{equation}
que coincide com a expressão $\sqrt{K}\cdot e/\Delta$ da análise de agregação do
Capítulo~4 da dissertação.

\paragraph{Codificação canônica CKKS.}
O codificador implementa o mapa de imersão canônica $\sigma$~\cite{Cheon2017}: um vetor
de $N/2$ valores complexos (os \emph{slots}) é associado a um polinômio de $\Rq$ por
avaliação nas $N$ raízes primitivas $2N$-ésimas $\omega_j = e^{i\pi(2j+1)/N}$. A
simetria hermitiana $\omega_{N-1-j} = \overline{\omega_j}$ garante que os coeficientes
do polinômio resultante sejam reais. Na codificação ($\sigma^{-1}$), o vetor de slots é
simetrizado, submetido a uma FFT de Cooley--Tukey radix-2 in-place, torcido por
$e^{-i\pi k/N}$, multiplicado pela escala $\Delta$ e arredondado para inteiros em
$[0,q)$. Na decodificação ($\sigma$), os coeficientes são centralizados em
$(-q/2,\,q/2]$, divididos por $\Delta$, torcidos por $e^{i\pi k/N}$ e submetidos à
IFFT. O erro de quantização por \emph{slot} é $\sim N/(2\Delta)$.

% =============================================================================
\section{Serialização, Compilação e Síntese}
\label{sec:serial_build}
% =============================================================================

\paragraph{Serialização.}
O módulo de serialização adota um formato binário \emph{little-endian} auto-descritivo.
A unidade atômica é a representação de um polinômio:
$[\,N\text{: u32, 4\,B}\,]\;[\,q\text{: u64, 8\,B}\,]\;[\,N \times\text{u64, }8N\text{\,B}\,]$,
totalizando $12 + 8N$ bytes. Um \texttt{Ciphertext} ou uma \texttt{PublicKey} consistem
em dois desses blocos ($2(12+8N)$ bytes); uma \texttt{SecretKey} consiste em um único
bloco ($12+8N$ bytes). Cada polinômio carrega seu próprio $N$ e $q$, permitindo
validação completa de integridade na desserialização: tamanho do \emph{buffer},
concordância de $(N,q)$ e ausência de bytes residuais.

\paragraph{Compilação e vetorização.}
O núcleo \texttt{ckks\_core} é compilado como biblioteca estática com código independente
de posição (\texttt{-fPIC}). As diretivas de compilação são: \texttt{-O3} (otimizações
agressivas), \texttt{-march=native} (instruções AVX2), \texttt{-ffast-math}
(otimizações de ponto flutuante, utilizadas principalmente pelo codificador) e
\texttt{-fopenmp-simd} (ativa \texttt{\#pragma omp simd} sem o \emph{runtime} de
\emph{threads} OpenMP, reservado para paralelização futura da NTT). Existe uma
divergência entre os dois caminhos de construção: o \texttt{setup.py} (modo
\emph{pip}) usa apenas \texttt{-O3 -march=native -Wall}, omitindo
\texttt{-ffast-math} e \texttt{-fopenmp-simd}.

\paragraph{Síntese e trabalhos futuros.}
Os três pilares técnicos do \texttt{ckks\_core} são: (i)~redução de Montgomery, que
elimina divisões do caminho crítico; (ii)~NTT negacíclica $O(N\log N)$ com tabelas de
twiddle pré-computadas em ordem bit-reversa e em forma de Montgomery; e (iii)~aritmética
modular sem desvios condicionais, que viabiliza a vetorização AVX2. As simplificações
de projeto (primo único, regime somente-aditivo) mantêm a implementação concisa sem
comprometer a correção para o caso de uso alvo: agregação federada de gradientes por
adição homomórfica. Os trabalhos futuros imediatos incluem a paralelização por
\emph{threads} dos laços da NTT, a substituição da multiplicação escolar $O(N^2)$ de
\texttt{mul\_plain} por \texttt{poly\_ntt\_mul}, e a extensão para cadeia RNS com
\emph{rescaling}, habilitando profundidade multiplicativa para circuitos de FL mais
complexos.

% =============================================================================
\begin{thebibliography}{9}

\bibitem{Cheon2017}
J.~H. Cheon, A.~Kim, M.~Kim, e Y.~Song,
``Homomorphic encryption for arithmetic of approximate numbers,''
em \emph{ASIACRYPT 2017}, LNCS vol.~10624, Springer, pp.~409--437.

\bibitem{Montgomery1985}
P.~L. Montgomery,
``Modular multiplication without trial division,''
\emph{Mathematics of Computation}, vol.~44, n.~170, pp.~519--521, 1985.

\bibitem{LongaNaehrig2016}
P.~Longa e M.~Naehrig,
``Speeding up the number theoretic transform for faster ideal lattice-based cryptography,''
em \emph{CANS 2016}, LNCS vol.~10052, Springer, pp.~124--139.

\bibitem{LPR2010}
V.~Lyubashevsky, C.~Peikert, e O.~Regev,
``On ideal lattices and learning with errors over rings,''
em \emph{EUROCRYPT 2010}, LNCS vol.~6110, Springer, pp.~1--23.

\bibitem{sealmanual}
Microsoft Research,
\emph{SEAL: Simple Encrypted Arithmetic Library -- Manual}, v.~4.1, 2023.
\url{https://github.com/microsoft/SEAL}

\bibitem{Pomerance1980}
C.~Pomerance, J.~L. Selfridge, e S.~S. Wagstaff~Jr.,
``The pseudoprimes to $25\cdot10^9$,''
\emph{Mathematics of Computation}, vol.~35, n.~151, pp.~1003--1026, 1980.

\end{thebibliography}

\end{document}
